\begin{abstract}
La recuperaci\'on aumentada por generaci\'on (RAG) permite responder sobre un corpus sin reentrenar el modelo, pero su calidad depende de decisiones que rara vez se documentan juntas: limpieza, segmentaci\'on, codificador, \'indice, reordenamiento, prompt y control de abstenci\'on. Este trabajo presenta un asistente RAG reproducible sobre 169.248 documentos recolectados de prensa colombiana e internacional y Bluesky, 161.868 indexados tras descartar p\'aginas de error y 618.286 fragmentos en Milvus, con modelos locales. La comparaci\'on de capas es exploratoria sobre el gold completo. Se justifica cada capa con experimentos pareados: la limpieza conservadora no mejora de forma significativa, el chunking fijo tiende a superar a la segmentaci\'on por oraciones y es m\'as r\'apido, bge-m3 supera a e5-base y a la fusi\'on con BM25, y un cross-encoder sobre el top-5 sube el nDCG de 0,809 a 0,865 y el P@1 de 0,737 a 0,868 sin perder cobertura. En el gold exploratorio el sistema entrega 37 de 38 preguntas respondibles, se abstiene en las 5 sin respuesta y acierta 25 de 37 de forma exacta. En un conjunto reservado de 20 preguntas nuevas, congelado antes de ejecutar, entrega 17, acierta 15 de forma exacta y se abstiene en las 5 sin respuesta; el recall@10 baja a 0,600 y la compuerta rechaza las 5 sin respuesta y 3 de las 20 directas. En conversaci\'on alcanza 14 de 14 con compuerta de dominio. RAGAS reporta 0,952 de fidelidad, 0,881 de precisi\'on de contexto, 0,921 de cobertura y 0,657 de correcci\'on. Se documentan tambi\'en las mejoras rechazadas, el coste de cada decisi\'on y los l\'imites del corpus.
\end{abstract}

\begin{IEEEkeywords}
recuperaci\'on aumentada por generaci\'on, Milvus, bge-m3, reordenamiento, evaluaci\'on de sistemas de recuperaci\'on, compuerta de dominio.
\end{IEEEkeywords}
